\chapter{结论}

本文以二甲醚（DME）为燃料，采用详细化学反应动力学机理和HCT化学动力学计算程序，对DME均质压燃燃烧过程进行了数值模拟研究。主要结论如下：

\begin{enumerate}
	\item DME的HCCI燃烧过程由低温反应放热阶段和高温反应放热阶段组成。低温反应阶段主要发生DME的低温氧化反应，产生大量中间产物；高温反应阶段发生快速燃烧放热，完成主要燃烧过程。

	\item 压缩比、当量比和进气温度的增加均使着火时刻提前。压缩比从8增加到18时，着火时刻提前约15°CA；当量比从0.2增加到0.5时，着火时刻提前约4°CA；进气温度从350K增加到450K时，着火时刻提前约10°CA。

	\item 发动机转速的增加使着火时刻推迟，这是因为转速增加导致反应时间缩短。

	\item 采用冷EGR可有效推迟着火时刻，EGR率从0\%增加到40\%时，着火时刻推迟约8°CA。EGR主要通过稀释效应、热效应和化学效应影响着火时刻。

	\item H2O2、H2和CO的添加使着火提前，而CH4和CH3OH的添加使着火推迟。其中H2O2对着火时刻的影响最为显著。
\end{enumerate}

本文的研究工作为DME均质压燃燃烧的优化控制提供了理论依据，对HCCI发动机的实用化具有一定的参考价值。未来的研究工作可以进一步考虑多维CFD模型与详细化学动力学模型的耦合，以更准确地预测HCCI燃烧过程。
